\section{Completed integral averages on the active mixed space}
\label{sec:mpo_mixed_wha_parent}

The weak-Hopf construction of Appendix B of
\cite{GarreRubio2023MPOSym} applies to the active physical space of the
mixed interpolation. The literal enlarged tensor has an obstruction to a
unital realization. The conditional construction below distinguishes
averaging from the two ambient-complement penalties; the realization
from fusion data remains a separate hypothesis
\cite{gap:glm23_mixed_weak_hopf_realization}.

\begin{lemma}[Zero physical summands obstruct unital realization]
    \label{lem:mpo_mixed_wha_obstruction}
    \leanok
    \lean{MPOTensor.WeakHopfRepresentation.tensor_diag_ne_zero,
        MPSTensor.MPOSymmetry.commonFixedPointInclusion_last_row,
        MPSTensor.MPOSymmetry.commonFixedPointInclusion_middle_row,
        MPSTensor.MPOSymmetry.changePhysicalBasis_eq_zero_of_row,
        MPSTensor.MPOSymmetry.arbitraryPhysicalMixedMPO_ne_weakHopfTensor}
    \uses{def:glm_dual_representation_tensor,def:mpo_arbitrary_mixed_operator}
    If the physical representation $\phi$ is unital and the dual
    representation $\psi$ is faithful, every diagonal tensor letter
    $\psi(f_{ii})$ is nonzero. Consequently the enlarged common mixed
    tensor cannot equal such a tensor whenever $d_0+d_1>0$.
\end{lemma}
\begin{proof}
    \leanok
    Faithfulness and $\psi(f_{ii})=0$ would give $f_{ii}=0$, contradicting
    $f_{ii}(1)=\phi(1)_{ii}=1$. The enlarged common tensor is zero on
    both unused physical summands.
\end{proof}

\begin{definition}[Completed mixed integral interaction]
    \label{def:mpo_mixed_wha_completion}
    \leanok
    \lean{MPOTensor.WeakHopfRepresentation.averageMatrix,
        MPSTensor.MPOSymmetry.mixedWeakHopfCompletedInteraction,
        MPSTensor.MPOSymmetry.arbitraryPhysicalMixedWeakHopfParent}
    \uses{def:mpo_arbitrary_mixed_operator,thm:glm_wha_parent_average}
    On the active alphabet put $Q=\rho_2(1)$ and
    $E(B)=\sum\rho_2(\Omega_{(1)})B\rho_2(S_\Omega(\Omega_{(2)}))$.
    For the extended mixed interaction $h_\gamma$, put
    $\widehat h_\gamma=E(h_\gamma)+I-Q$.
    With $J$ the common physical inclusion and $J_2=J\otimes J$, put
    $\widetilde h_\gamma=I-J_2(I-\widehat h_\gamma)J_2^*$.
\end{definition}

\begin{lemma}[Continuous completion of the integral average]
    \label{lem:mpo_mixed_wha_completion}
    \leanok
    \lean{MPOTensor.WeakHopfRepresentation.averageMatrix_eq_sum,
        MPOTensor.WeakHopfRepresentation.continuous_averageMatrix,
        MPSTensor.MPOSymmetry.averageMatrix_mixedEndpointInteraction_of_mem_Ioo,
        MPSTensor.MPOSymmetry.averageMatrix_mixedEndpointInteraction,
        MPSTensor.MPOSymmetry.mixedWeakHopfCompletedInteraction_eq,
        MPSTensor.MPOSymmetry.continuous_mixedWeakHopfCompletedInteraction,
        MPSTensor.MPOSymmetry.mixedWeakHopfCompletedInteraction_posSemidef,
        MPSTensor.MPOSymmetry.arbitraryPhysicalMixedWeakHopfParent_eq,
        MPSTensor.MPOSymmetry.continuousOn_arbitraryPhysicalMixedWeakHopfParent,
        MPSTensor.MPOSymmetry.arbitraryPhysicalMixedWeakHopfParent_posSemidef}
    \uses{def:mpo_mixed_wha_completion,thm:glm_wha_parent_average}
    Assume injective endpoints of positive bond dimensions, exact
    biorthogonal action decompositions with a common multiplicity, and
    a normalized weak-Hopf realization of the actual active common
    operator. Then $E(h_\gamma)=Q-I+h_\gamma$ and
    $\widehat h_\gamma=h_\gamma$ throughout $[0,1]$.
    The enlarged completed terms equal the existing enlarged mixed
    interaction, are positive, and vary continuously on $[0,1]$.
\end{lemma}
\begin{proof}
    \leanok
    \uses{thm:glm_wha_parent_average}
    Interior injectivity and compatibility give $E(P_\gamma)=P_\gamma$.
    Since $E(I)=Q$, subtraction gives the stated average identity.
    Finite sums of fixed matrix multiplications are continuous, so
    density extends the identity to both endpoints. Adding $I-Q$ and
    then the unused physical-space penalty gives the claimed equality.
\end{proof}

\begin{lemma}[Symmetry and gap of the completed path]
    \label{lem:mpo_mixed_wha_symmetry_gap}
    \leanok
    \lean{MPOTensor.WeakHopfRepresentation.rep_one_commute_mpoWithBoundary,
        MPSTensor.MPOSymmetry.averageMatrix_mixedEndpointInteraction_commute,
        MPSTensor.MPOSymmetry.interactionHamiltonian_averageMatrix_mixedEndpointInteraction_commute,
        MPOTensor.WeakHopfRepresentation.exists_adjointBoundary_changePhysicalBasis,
        MPSTensor.MPOSymmetry.mixedWeakHopfCompletedInteraction_commute,
        MPSTensor.MPOSymmetry.interactionHamiltonian_mixedWeakHopfCompletedInteraction_commute,
        MPSTensor.MPOSymmetry.interactionHamiltonian_arbitraryPhysicalMixedWeakHopfParent_commute,
        MPSTensor.MPOSymmetry.exists_uniform_arbitraryPhysicalMixedWeakHopfParent_gap}
    \uses{lem:mpo_mixed_wha_completion,thm:mpo_arbitrary_parent_commutation_adjoint_closed}
    Under the preceding hypotheses, both active and enlarged completed
    periodic Hamiltonians commute with their common periodic operators
    for all $N\ge2$. For the enlarged Hamiltonians there is
    $0<\delta\le1$, independent of $\gamma\in[0,1]$ and $N\ge2$,
    such that $\delta\|v\|\le\|H_\gamma^N v\|$ for every vector
    orthogonal to the kernel. No separate adjoint-closure hypothesis is
    imposed, and no gap claim is made for the uncompleted average.
    If the weak-Hopf algebra is finite dimensional, the literal active
    average also commutes with every two-site boundary operator, and its
    periodic sum commutes with the common periodic operator for all $N\ge2$.
\end{lemma}
\begin{proof}
    \leanok
    \uses{lem:mpo_mixed_wha_completion,thm:mpo_arbitrary_parent_commutation_adjoint_closed}
    The realization supplies adjoint boundaries on the active space.
    For $J_N=J^{\otimes N}$, physical inclusion gives
    \[
        O^N_{\widetilde T,X}=J_N O^N_{T,X}J_N^*.
    \]
    Thus $(O^N_{T,X})^*=O^N_{T,Y}$ implies
    $(O^N_{\widetilde T,X})^*=O^N_{\widetilde T,Y}$, preserving the same
    adjoint boundary.
    Compatibility gives local commutation in the interior; continuity
    extends it to the endpoints, and cyclic placement gives the periodic
    assertion. Equality with the existing mixed path transfers its
    uniform gap without a spectral comparison of different Hamiltonians.
    In finite dimension, for $n>0$ every boundary operator has the form
    $B=O^n_{T,X}=\rho_n(x)$. Multiplicativity gives
    \[
        \rho_n(1)\rho_n(x)=\rho_n(1x)=\rho_n(x1)
        =\rho_n(x)\rho_n(1).
    \]
    In particular, at $n=2$ we have $\rho_2(1)=Q$ and $QB=BQ$;
    adding $Q-I$ therefore preserves commutation with $B$.
\end{proof}
